% =========================================================
% template.tex  (NO tiene \documentclass ni \begin{document})
    % Plantilla institucional robusta: modo claro,
    % fallbacks de imágenes, minted, entornos de código claros/oscuros y macros docentes.
    % =========================================================

    % -------------------------
    % CONFIG POR DEFECTO
    % (En main.tex puedes sobrescribir con \renewcommand)
    % -------------------------
    \newcommand{\CourseName}{Laboratorio Inteligencia Artificial}
    \newcommand{\CourseCode}{ING01336}
    \newcommand{\Faculty}{Facultad de Ingenierías}
    \newcommand{\Institution}{Politécnico Colombiano Jaime Isaza Cadavid}
    \newcommand{\Instructor}{Deimer Miranda Montoya, MSc.(c).}
    \newcommand{\Email}{deimer\_miranda91162@elpoli.edu.co}
    \newcommand{\InstitutionLogo}{Template/LogoSena.png}

    % Control global de imágenes: 1 = mostrar, 0 = no cargar
    \newcommand{\UseImages}{1}

    % -------------------------
    % MOTOR Y FUENTES (ROBUSTO)
    % -------------------------
    \usepackage{iftex}
    \ifPDFTeX
    \usepackage[utf8]{inputenc}
    \usepackage[T1]{fontenc}
    \usepackage{lmodern}
    \else
    \usepackage{fontspec}
    \usepackage[sfdefault]{noto}
    \fi



    % -------------------------
    % PAQUETES
    % -------------------------
    \usetheme{Madrid}
    \usepackage{xparse}
    \usepackage{fontawesome5}
    \usepackage{booktabs}
    \usepackage{xcolor}
    \usepackage{graphicx}
    \usepackage{tikz}
    \usetikzlibrary{calc,positioning,arrows.meta,mindmap,shapes.geometric}
    \usepackage{hyperref}
    \usepackage{tcolorbox}
    \tcbuselibrary{minted,skins,breakable}

    \hypersetup{colorlinks=true, linkcolor=blue, urlcolor=blue, citecolor=blue}

    % -------------------------
    % COLORES (INSTITUCIONAL)
    % -------------------------
    \definecolor{BrandYellow}{RGB}{255,185,62}
    \definecolor{BrandGreen}{RGB}{20,196,134}
    \definecolor{BrandBlue}{RGB}{30,60,120}
    \definecolor{DarkBG}{RGB}{14,16,20}
    \definecolor{DarkFG}{RGB}{235,235,235}

    \definecolor{SoftGray}{RGB}{245,246,248}
    \definecolor{MediumGray}{RGB}{210,213,218}
    \definecolor{DarkGray}{RGB}{45,47,52}

    % Fondos de código
    \definecolor{pybg}{rgb}{0.12,0.18,0.30}
    \definecolor{cppbg}{rgb}{0.10,0.10,0.10}
    \definecolor{bashbg}{rgb}{0.2,0.2,0.2}

    % Fondos específicos para código HTML
    \definecolor{HtmlClaroFondo}{HTML}{F7F7F7}
    \definecolor{HtmlClaroBorde}{HTML}{D0D0D0}
    \definecolor{HtmlOscuroFondo}{HTML}{1E1E1E}
    \definecolor{HtmlOscuroBorde}{HTML}{444444}

    % -------------------------
    % COLORES PARA BLOQUES (DIFERENTES)
    % -------------------------
    % Idea
    \definecolor{IdeaTitle}{RGB}{30,60,120}     % azul
    \definecolor{IdeaBody}{RGB}{235,241,255}    % azul muy claro

    % Practice
    \definecolor{PracticeTitle}{RGB}{20,120,90} % verde/teal
    \definecolor{PracticeBody}{RGB}{232,250,245}

    % Definition
    \definecolor{DefTitle}{RGB}{90,70,160}      % violeta
    \definecolor{DefBody}{RGB}{240,236,255}

    % Result
    \definecolor{ResultTitle}{RGB}{0,120,0}     % verde
    \definecolor{ResultBody}{RGB}{235,250,235}

    % Warning (alert)
    \definecolor{WarnTitle}{RGB}{160,30,30}     % rojo
    \definecolor{WarnBody}{RGB}{255,235,235}

    % Comparación
    \definecolor{CompareTitle}{RGB}{28,92,145}
    \definecolor{CompareBody}{RGB}{235,245,252}

    % Pregunta
    \definecolor{QuestionTitle}{RGB}{176,112,20}
    \definecolor{QuestionBody}{RGB}{255,247,229}

    % Ejemplo
    \definecolor{ExampleTitle}{RGB}{84,78,150}
    \definecolor{ExampleBody}{RGB}{243,241,255}

    % =========================================================
    % TEMA VISUAL - SOLO MODO CLARO
    % =========================================================
    \setbeamercolor{background canvas}{bg=white}
    \setbeamercolor{normal text}{fg=black,bg=white}
    \setbeamercolor{structure}{fg=BrandBlue}
    \setbeamercolor{itemize item}{fg=BrandBlue}
    \setbeamercolor{itemize subitem}{fg=BrandBlue}
    \setbeamercolor{enumerate item}{fg=BrandBlue}
    \setbeamercolor{alerted text}{fg=WarnTitle}
    \setbeamercolor{example text}{fg=PracticeTitle}
    \setbeamercolor{caption name}{fg=BrandBlue}

    % -------------------------
    % Fondo tipo banda superior (frametitle)
    % -------------------------
    \setbeamercolor{frametitle}{fg=white,bg=BrandBlue}
    \setbeamertemplate{frametitle}{
        \nointerlineskip
        \begin{beamercolorbox}[wd=\paperwidth,ht=0.5cm,dp=0.2cm,leftskip=0.6cm]{frametitle}
            \usebeamerfont{frametitle}\insertframetitle
        \end{beamercolorbox}
    }

    % -------------------------
    % IMÁGENES (ROBUSTAS)
    % -------------------------
    \graphicspath{{Template/}{Images/}{Images2/}{Logos/}}

    \newcommand{\smartimage}[2][0.9\linewidth]{%
        \ifnum\UseImages=1
        \IfFileExists{#2}{%
            \includegraphics[width=#1]{#2}%
        }{%
            \begingroup
            \setlength{\fboxsep}{4pt}%
            \fbox{\parbox{#1}{\centering\scriptsize Imagen no encontrada:\\ \texttt{#2}}}%
            \endgroup
        }%
        \fi
    }

    \newcommand{\smartlogo}[2][1.8cm]{%
        \ifnum\UseImages=1
        \IfFileExists{#2}{%
            \includegraphics[height=#1]{#2}%
        }{%
            {\footnotesize\textbf{[Logo no encontrado]}}%
        }%
        \fi
    }

    % -------------------------
    % CÓDIGO (MINTED)
    % Se conservan entornos claros y oscuros de código independientemente del tema visual.
    % Requiere compilar con -shell-escape y pygments instalado
    % -------------------------
    \usepackage{minted}
    \usemintedstyle{monokai}

    % =========================================================
    % COLORES PARA ENTORNOS DE CÓDIGO
    % =========================================================

    % Fondo claro general
    \definecolor{CodigoClaroFondo}{HTML}{F7F7F7}
    \definecolor{CodigoClaroBorde}{HTML}{D0D0D0}

    % Fondo oscuro general
    \definecolor{CodigoOscuroFondo}{HTML}{1E1E1E}
    \definecolor{CodigoOscuroBorde}{HTML}{444444}

    % =========================================================
    % ENTORNOS HTML CON RESALTADO DE SINTAXIS
    % Todos los frames que los utilicen deben llevar [fragile].
    % =========================================================

    % Fragmentos cortos HTML con fondo claro y sin numeración
    \newtcblisting{htmlcode}{
        listing engine=minted,
        minted language=html,
        minted style=friendly,
        minted options={
            fontsize=\small,
            breaklines=true,
            autogobble=true,
            stripnl=true,
            tabsize=4,
            encoding=utf8
        },
        colback=HtmlClaroFondo,
        colframe=HtmlClaroBorde,
        coltext=black,
        boxrule=0.6pt,
        arc=2mm,
        left=2mm,
        right=2mm,
        top=0.0mm,
        bottom=0.0mm,
        before skip=1pt,
        after skip=1pt,
        listing only,
        enhanced
    }

    % Fragmentos cortos HTML con fondo oscuro y sin numeración
    \newtcblisting{htmlcodedark}{
        listing engine=minted,
        minted language=html,
        minted style=monokai,
        minted options={
            fontsize=\small,
            breaklines=true,
            autogobble=true,
            stripnl=true,
            tabsize=4,
            encoding=utf8
        },
        colback=HtmlOscuroFondo,
        colframe=HtmlOscuroBorde,
        coltext=white,
        boxrule=0.6pt,
        arc=2mm,
        left=2mm,
        right=2mm,
        top=0.0mm,
        bottom=0.0mm,
        before skip=1pt,
        after skip=1pt,
        listing only,
        enhanced
    }

    % =========================================================
    % PYTHON
    % =========================================================

    % Python con fondo claro
    \newtcblisting{pythoncode}{
        listing engine=minted,
        minted language=python,
        minted style=friendly,
        minted options={
            fontsize=\small,
            breaklines=true,
            autogobble=true,
            stripnl=true,
            tabsize=4,
            encoding=utf8
        },
        colback=CodigoClaroFondo,
        colframe=CodigoClaroBorde,
        coltext=black,
        boxrule=0.6pt,
        arc=2mm,
        left=2mm,
        right=2mm,
        top=0.0mm,
        bottom=0.0mm,
        before skip=1pt,
        after skip=1pt,
        listing only,
        enhanced
    }

    % Python con fondo oscuro
    \newtcblisting{pythoncodedark}{
        listing engine=minted,
        minted language=python,
        minted style=monokai,
        minted options={
            fontsize=\small,
            breaklines=true,
            autogobble=true,
            stripnl=true,
            tabsize=4,
            encoding=utf8
        },
        colback=CodigoOscuroFondo,
        colframe=CodigoOscuroBorde,
        coltext=white,
        boxrule=0.6pt,
        arc=2mm,
        left=2mm,
        right=2mm,
        top=0.0mm,
        bottom=0.0mm,
        before skip=1pt,
        after skip=1pt,
        listing only,
        enhanced
    }

    % =========================================================
    % C++
    % =========================================================

    % C++ con fondo claro
    \newtcblisting{cppcode}{
        listing engine=minted,
        minted language=cpp,
        minted style=friendly,
        minted options={
            fontsize=\small,
            breaklines=true,
            autogobble=true,
            stripnl=true,
            tabsize=4,
            encoding=utf8
        },
        colback=CodigoClaroFondo,
        colframe=CodigoClaroBorde,
        coltext=black,
        boxrule=0.6pt,
        arc=2mm,
        left=2mm,
        right=2mm,
        top=0.0mm,
        bottom=0.0mm,
        before skip=1pt,
        after skip=1pt,
        listing only,
        enhanced
    }

    % C++ con fondo oscuro
    \newtcblisting{cppcodedark}{
        listing engine=minted,
        minted language=cpp,
        minted style=monokai,
        minted options={
            fontsize=\small,
            breaklines=true,
            autogobble=true,
            stripnl=true,
            tabsize=4,
            encoding=utf8
        },
        colback=CodigoOscuroFondo,
        colframe=CodigoOscuroBorde,
        coltext=white,
        boxrule=0.6pt,
        arc=2mm,
        left=2mm,
        right=2mm,
        top=0.0mm,
        bottom=0.0mm,
        before skip=1pt,
        after skip=1pt,
        listing only,
        enhanced
    }

    % =========================================================
    % CSS
    % =========================================================

    % CSS con fondo claro
    \newtcblisting{csscode}{
        listing engine=minted,
        minted language=css,
        minted style=friendly,
        minted options={
            fontsize=\small,
            breaklines=true,
            autogobble=true,
            stripnl=true,
            tabsize=4,
            encoding=utf8
        },
        colback=CodigoClaroFondo,
        colframe=CodigoClaroBorde,
        coltext=black,
        boxrule=0.6pt,
        arc=2mm,
        left=2mm,
        right=2mm,
        top=0.0mm,
        bottom=0.0mm,
        listing only,
        enhanced
    }

    % CSS con fondo oscuro
    \newtcblisting{csscodedark}{
        listing engine=minted,
        minted language=css,
        minted style=monokai,
        minted options={
            fontsize=\small,
            breaklines=true,
            autogobble=true,
            stripnl=true,
            tabsize=4,
            encoding=utf8
        },
        colback=CodigoOscuroFondo,
        colframe=CodigoOscuroBorde,
        coltext=white,
        boxrule=0.6pt,
        arc=2mm,
        left=2mm,
        right=2mm,
        top=0.0mm,
        bottom=0.0mm,
        listing only,
        enhanced
    }

    % =========================================================
    % SQL
    % =========================================================

    % SQL con fondo claro
    \newtcblisting{sqlcode}{
        listing engine=minted,
        minted language=sql,
        minted style=friendly,
        minted options={
            fontsize=\small,
            breaklines=true,
            autogobble=true,
            stripnl=true,
            tabsize=2,
            encoding=utf8
        },
        colback=CodigoClaroFondo,
        colframe=CodigoClaroBorde,
        coltext=black,
        boxrule=0.6pt,
        arc=2mm,
        left=2mm,
        right=2mm,
        top=0.0mm,
%        before skip=1pt,
%        after skip=1pt,
        bottom=0.0mm,
        listing only,
        enhanced
    }

    % SQL con fondo oscuro
    \newtcblisting{sqlcodedark}{
        listing engine=minted,
        minted language=sql,
        minted style=monokai,
        minted options={
            fontsize=\small,
            breaklines=true,
            autogobble=true,
            stripnl=true,
            tabsize=2,
            encoding=utf8
        },
        colback=CodigoOscuroFondo,
        colframe=CodigoOscuroBorde,
        coltext=white,
        boxrule=0.6pt,
        arc=2mm,
        left=2mm,
        right=2mm,
        top=0.0mm,
        bottom=0.0mm,
%        before skip=1pt,
%        after skip=1pt,
        listing only,
        enhanced
    }

    % =========================================================
    % JAVA
    % =========================================================

    % Java con fondo claro
    \newtcblisting{javacode}{
        listing engine=minted,
        minted language=java,
        minted style=friendly,
        minted options={
            fontsize=\small,
            breaklines=true,
            autogobble=true,
            stripnl=true,
            tabsize=4,
            encoding=utf8
        },
        colback=CodigoClaroFondo,
        colframe=CodigoClaroBorde,
        coltext=black,
        boxrule=0.6pt,
        arc=2mm,
        left=2mm,
        right=2mm,
        top=0.0mm,
        bottom=0.0mm,
%        before skip=1pt,
%        after skip=1pt,
        listing only,
        enhanced
    }

    % Java con fondo oscuro
    \newtcblisting{javacodedark}{
        listing engine=minted,
        minted language=java,
        minted style=monokai,
        minted options={
            fontsize=\small,
            breaklines=true,
            autogobble=true,
            stripnl=true,
            tabsize=4,
            encoding=utf8
        },
        colback=CodigoOscuroFondo,
        colframe=CodigoOscuroBorde,
        coltext=white,
        boxrule=0.6pt,
        arc=2mm,
        left=2mm,
        right=2mm,
        top=0.0mm,
        bottom=0.0mm,
%        before skip=1pt,
%        after skip=1pt,
        listing only,
        enhanced
    }

    % =========================================================
    % BASH
    % =========================================================

    % Bash con fondo claro
    \newtcblisting{bashcode}{
        listing engine=minted,
        minted language=bash,
        minted style=friendly,
        minted options={
            fontsize=\small,
            breaklines=true,
            autogobble=true,
            stripnl=true,
            tabsize=4,
            encoding=utf8
        },
        colback=CodigoClaroFondo,
        colframe=CodigoClaroBorde,
        coltext=black,
        boxrule=0.6pt,
        arc=2mm,
        left=2mm,
        right=2mm,
        top=0.0mm,
        bottom=0.0mm,
%        before skip=1pt,
%        after skip=1pt,
        listing only,
        enhanced
    }

    % Bash con fondo oscuro
    \newtcblisting{bashcodedark}{
        listing engine=minted,
        minted language=bash,
        minted style=monokai,
        minted options={
            fontsize=\small,
            breaklines=true,
            autogobble=true,
            stripnl=true,
            tabsize=4,
            encoding=utf8
        },
        colback=CodigoOscuroFondo,
        colframe=CodigoOscuroBorde,
        coltext=white,
        boxrule=0.6pt,
        arc=2mm,
        left=2mm,
        right=2mm,
        top=0.0mm,
        bottom=0.0mm,
%        before skip=1pt,
%        after skip=1pt,
        listing only,
        enhanced
    }

    % =========================================================
    % MONGODB / MONGOSH - FONDO CLARO
    % =========================================================

    \newtcblisting{mongodbcode}{
        listing engine=minted,
        minted language=javascript,
        minted style=friendly,
        minted options={
            fontsize=\small,
            breaklines=true,
            autogobble=true,
            stripnl=true,
            tabsize=4,
            encoding=utf8
        },
        colback=CodigoClaroFondo,
        colframe=CodigoClaroBorde,
        coltext=black,
        boxrule=0.6pt,
        arc=2mm,
        left=2mm,
        right=2mm,
        top=0.0mm,
        bottom=0.0mm,
        listing only,
        enhanced,
        breakable
    }

    % =========================================================
    % MONGODB / MONGOSH - FONDO OSCURO
    % =========================================================

    \newtcblisting{mongodbcodedark}{
        listing engine=minted,
        minted language=javascript,
        minted style=monokai,
        minted options={
            fontsize=\small,
            breaklines=true,
            autogobble=true,
            stripnl=true,
            tabsize=4,
            encoding=utf8
        },
        colback=CodigoOscuroFondo,
        colframe=CodigoOscuroBorde,
        coltext=white,
        boxrule=0.6pt,
        arc=2mm,
        left=2mm,
        right=2mm,
        top=0.0mm,
        bottom=0.0mm,
        listing only,
        enhanced,
        breakable
    }

    % =========================================================
    % CONSOLA / SALIDA DE TERMINAL - TEMA OSCURO
    % =========================================================
    \newtcblisting{consolecodedark}{
        listing engine=minted,
        minted language=text,
        minted options={
            breaklines,
            stripnl=true,
            breakanywhere,
            fontsize=\small,
            autogobble
        },
        listing only,
        enhanced,
        breakable,
        colback=black!92,
        colframe=black!65,
        coltext=white,
        boxrule=0.6pt,
        arc=2mm,
        left=2mm,
        right=2mm,
        top=0.0mm,
        bottom=0.0mm,
        before skip=2mm,
        after skip=2mm
    }
    % =========================================================
    % CONSOLA / SALIDA DE TERMINAL - TEMA CLARO
    % =========================================================
    \newtcblisting{consolecode}{
        listing engine=minted,
        minted language=text,
        minted options={
            breaklines,
            stripnl=true,
            breakanywhere,
            fontsize=\small,
            autogobble
        },
        listing only,
        enhanced,
        breakable,
        colback=black!4,
        colframe=black!30,
        coltext=black,
        boxrule=0.6pt,
        arc=2mm,
        left=2mm,
        right=2mm,
        top=0.0mm,
        bottom=0.0mm,
        before skip=2mm,
        after skip=2mm
    }

    % -------------------------
    % MACROS DOCENTES
    % -------------------------
    \newcommand{\SectionSlide}[2]{%
        \begin{frame}
            \vspace*{0.1cm}
            \centering
            \begin{tikzpicture}
                \node[fill=BrandBlue, rounded corners,
                minimum width=0.86\textwidth, minimum height=1.6cm,
                text width=0.82\textwidth, align=center] (box)
                {\color{white}\textbf{\huge #1}\\[2pt]\Large #2};
            \end{tikzpicture}
        \end{frame}
    }

    % =========================================================
    % BLOQUES DOCENTES CON TÍTULO OPCIONAL + COLORES DIFERENTES
    % =========================================================
    % Idea: \Idea[Mi título]{Contenido}
    \NewDocumentCommand{\Idea}{O{Idea clave} +m}{
        \begingroup
        \setbeamercolor{block title}{fg=white,bg=IdeaTitle}
        \setbeamercolor{block body}{fg=black,bg=IdeaBody}
        \begin{block}{\faLightbulb\ \,#1}
            #2
        \end{block}
        \endgroup
    }

    % Practice
    \NewDocumentCommand{\Practice}{O{Actividad práctica} +m}{
        \begingroup
        \setbeamercolor{block title}{fg=white,bg=PracticeTitle}
        \setbeamercolor{block body}{fg=black,bg=PracticeBody}
        \begin{block}{\faLaptopCode\ \,#1}
            #2
        \end{block}
        \endgroup
    }

    % Definition
    \NewDocumentCommand{\DefBlock}{O{Definición} +m}{
        \begingroup
        \setbeamercolor{block title}{fg=white,bg=DefTitle}
        \setbeamercolor{block body}{fg=black,bg=DefBody}
        \begin{block}{\faBook\ \,#1}
            #2
        \end{block}
        \endgroup
    }

    % Result
    \NewDocumentCommand{\Result}{O{Resultado} +m}{
        \begingroup
        \setbeamercolor{block title}{fg=white,bg=ResultTitle}
        \setbeamercolor{block body}{fg=black,bg=ResultBody}
        \begin{block}{\faCheckCircle\ \,#1}
            #2
        \end{block}
        \endgroup
    }

    % Warning (alertblock)
    \NewDocumentCommand{\Warning}{O{Atención} +m}{
        \begingroup
        \setbeamercolor{block title alerted}{fg=white,bg=WarnTitle}
        \setbeamercolor{block body alerted}{fg=black,bg=WarnBody}
        \begin{alertblock}{\faExclamationTriangle\ \,#1}
            #2
        \end{alertblock}
        \endgroup
    }

    % Comparación: \Compare[Mi título]{Contenido}
    \NewDocumentCommand{\Compare}{O{Comparación} +m}{
        \begingroup
        \setbeamercolor{block title}{fg=white,bg=CompareTitle}
        \setbeamercolor{block body}{fg=black,bg=CompareBody}
        \begin{block}{\faBalanceScale\ \,#1}
            #2
        \end{block}
        \endgroup
    }

    % Pregunta: \Question[Mi título]{Contenido}
    \NewDocumentCommand{\Question}{O{Pregunta guía} +m}{
        \begingroup
        \setbeamercolor{block title}{fg=white,bg=QuestionTitle}
        \setbeamercolor{block body}{fg=black,bg=QuestionBody}
        \begin{block}{\faQuestionCircle\ \,#1}
            #2
        \end{block}
        \endgroup
    }

    % Ejemplo: \ExampleBlock[Mi título]{Contenido}
    % Se usa ExampleBlock para evitar colisión con comandos propios de Beamer.
    \NewDocumentCommand{\ExampleBlock}{O{Ejemplo} +m}{
        \begingroup
        \setbeamercolor{block title}{fg=white,bg=ExampleTitle}
        \setbeamercolor{block body}{fg=black,bg=ExampleBody}
        \begin{block}{\faPuzzlePiece\ \,#1}
            #2
        \end{block}
        \endgroup
    }

    \newcommand{\AgendaSlide}{%
        \begin{frame}{Agenda}
            \tableofcontents
        \end{frame}
    }

    % =========================================================
    % SLIDE DE NORMAS INSTITUCIONALES
    % =========================================================
    \newcommand{\NormasSlide}{%
        \begin{frame}{Normas del curso}
            \begin{minipage}{0.32\textwidth}
                \centering
                \smartimage[1cm]{Template/Puntual.png}\\[0.2cm]
                \textbf{Cumple con los horarios y demuestra compromiso académico en cada sesión.}
            \end{minipage}\hfill
            \begin{minipage}{0.32\textwidth}
                \centering
                \smartimage[1cm]{Template/Actitud.png}\\[0.2cm]
                \textbf{Mantén una actitud profesional, ética y orientada al aprendizaje continuo.}
            \end{minipage}\hfill
            \begin{minipage}{0.32\textwidth}
                \centering
                \smartimage[1cm]{Template/Comunicacion.png}\\[0.2cm]
                \textbf{Comunica tus ideas con claridad, respeto y fundamento técnico.}
            \end{minipage}

            \vspace{0.2cm}

            \begin{minipage}{0.48\textwidth}
                \centering
                \smartimage[1cm]{Template/Participacion.png}\\[0.2cm]
                \textbf{Participa activamente y contribuye al trabajo colaborativo.}
            \end{minipage}\hfill
            \begin{minipage}{0.48\textwidth}
                \centering
                \smartimage[1cm]{Template/Cuidado.png}\\[0.2cm]
                \textbf{Haz uso responsable de los recursos, equipos e instalaciones.}
            \end{minipage}
        \end{frame}
    }


    % -------------------------
    % PIE DE PÁGINA
    % -------------------------
    \setbeamertemplate{navigation symbols}{}

    \setbeamertemplate{footline}{
        \leavevmode
        \hbox to \paperwidth{%
            \hspace{0.35cm}%
            % Zona izquierda: navegación y numeración
            \makebox[0.22\paperwidth][l]{%
                {\tiny
                    \Acrobatmenu{PrevPage}{\textcolor{black}{$\triangleleft$}}%
                    \hspace{0.22cm}%
                    \insertpagenumber/\insertpresentationendpage%
                    \hspace{0.22cm}%
                    \Acrobatmenu{NextPage}{\textcolor{black}{$\triangleright$}}%
                }%
            }%
            \hfill%
            % Zona derecha: instructor + logo
            \makebox[0.70\paperwidth][r]{%
                {\tiny\Instructor}%
                \hspace{0.18cm}%
                \raisebox{-0.10cm}{\smartlogo[0.34cm]{\InstitutionLogo}}%
                %            \raisebox{-0.10cm}{\smartlogo[0.34cm]{Template/LogoSena.png}}%
            }%
            \hspace{0.35cm}%
        }%
        \vspace{0.08cm}
    }
    % -------------------------
    % PORTADA (UN COMANDO)
    % -------------------------
    \newcommand{\TitleSlide}{%
        \begin{frame}
            \centering
            %\smartlogo[1.8cm]{logo.png}

            \vspace{0.4cm}
            \color{BrandGreen}\rule{0.9\textwidth}{0.5pt}

            \vspace{0.6cm}
            {\color{black}\fontsize{24pt}{28pt}\selectfont \CourseName\\
                \CourseCode}

            \vfill
            {\color{black}\fontsize{16pt}{18pt}\selectfont \Faculty \par \Institution}

            \vfill
            {\color{black}\fontsize{14pt}{18pt}\selectfont \Instructor \par \Email}

            \vspace{0.6cm}
            \color{BrandGreen}\rule{0.9\textwidth}{0.5pt}
        \end{frame}
    }

% =========================================================
% ESPACIADO AUTOMÁTICO ANTES DE ECUACIONES DISPLAY
% =========================================================

\newlength{\DisplayMathTopSpace}
\setlength{\DisplayMathTopSpace}{-0.2cm}

\let\OriginalDisplayMath\[
\renewcommand{\[}{%
    \vspace{\DisplayMathTopSpace}%
    \OriginalDisplayMath
}
